\chapter{Mathematical Correspondence of the EasyCrypt Proof Objects}
\label{app:mathematical-correspondence}

\providecommand{\ecfile}[1]{\nolinkurl{#1}}
\providecommand{\ecsym}[1]{\nolinkurl{#1}}

Entries in this appendix use schematic metavariables when describing general
EasyCrypt syntax.  In particular, occurrences of ellipses such as
\ecsym{tau1 -> ... -> taun} denote families of language constructs, not omitted
parts of a specific HAETAE theorem.  Exact theorem and module names are written
without ellipses.

This appendix gives the mathematical reading of every kind of EasyCrypt object
used in the HAETAE provable-security development.  The goal is not to restate
the proof script tactically, but to say what the checked declarations mean as
mathematics.  Let \(\llbracket - \rrbracket\) denote the intended mathematical
interpretation of an EasyCrypt declaration in the manifest proof surface.

\section{Global interpretation principle}
\label{app:global-interpretation-principle}

The EasyCrypt development is a typed formalization of a game-based proof.
Mathematically, each manifest file contributes objects in a category of
discrete probabilistic programs, predicates over their states, and
inequalities between event probabilities.  The top-level checked implication is
therefore a fixed-ledger formal-model statement of the following form:
\[
  \mathcal{H}_{\mathsf{hop}}
  \wedge
  \mathcal{H}_{\mathsf{hard}}
  \Longrightarrow
  \Adv_{\mathsf{ECModel},\rom}^{\eufcma}(A)
  \leq
  B_{\mathsf{loss}}(2,2,\lambda)
  + L_{\mathsf{ModuleSIS}}
  + L_{\mathsf{MLWE}}
  + L_{\mathsf{bimodal}\to\mathsf{MSIS}},
\]
where the right-hand side is the concrete expression defined in the reduction
file after applying the declared hardness bounds.  The reductions \(B_i\) are
either explicit adapters in individual hop lemmas or abstract reduction modules
supplied through the final lemma's antecedents.  In particular, the final
composition lemma does not derive the MLWE or Module-SIS solvers from \(A\); it
uses the explicit hop/reduction and hardness premises present in the EasyCrypt
statement.  Each EasyCrypt type, operation, module, lemma, and Hoare/equivalence
judgment contributes one well-typed component to this implication.  A
specification-level or arbitrary-query HAETAE theorem would require additional
correspondence and parameterization work.

\section{Types as carrier sets}
\label{app:types-as-carrier-sets}

An EasyCrypt declaration of the form
\begin{center}
\ecsym{type t}
\end{center}
is interpreted as a carrier set \(\llbracket t \rrbracket\).  If the type is
abstract, the proof treats it as an arbitrary set satisfying only the
properties derivable from the surrounding declarations.  If the type is an
explicit finite datatype, such as the HAETAE mode type, it is interpreted as the
listed finite set.  Product, list, and distribution types are interpreted in the
usual set-theoretic way:
\[
  \llbracket \tau_1 * \tau_2 \rrbracket
  =
  \llbracket \tau_1 \rrbracket \times \llbracket \tau_2 \rrbracket,
  \qquad
  \llbracket \mathsf{List}(\tau) \rrbracket
  =
  \llbracket \tau \rrbracket^{*},
\]
and \(\llbracket \mathsf{Distr}(\tau) \rrbracket\) is the set of discrete
probability or subprobability laws over \(\llbracket \tau \rrbracket\), with
\ecsym{is\_lossless} asserting total mass \(1\).

Thus all manifest types have one of the following mathematical roles.

\begin{longtable}{p{0.25\textwidth}p{0.65\textwidth}}
\caption{Mathematical reading of all manifest type declarations.}
\label{tab:type-correspondence}\\
\toprule
EasyCrypt type class & Mathematical correspondence \\
\midrule
\endfirsthead
\toprule
EasyCrypt type class & Mathematical correspondence \\
\midrule
\endhead
Signature-game types &
Carrier sets for public keys, secret keys, messages, contexts, signatures,
random-oracle queries, random-oracle outputs, and signed-query records.  These
sets define the sample space of the generic correctness, EUF-CMA, SUF-CMA, and
UF-NMA experiments. \\
HAETAE algebra types &
Concrete carrier sets for byte strings, seeds, CRHs, random coins, coefficients,
polynomials, challenges, polynomial vectors, matrices, public keys, secret keys,
and signatures.  These sets are modeled as integers, lists, products, and
mode-indexed list families. \\
Parameter-mode types &
Finite sets of parameter choices.  The declaration of HAETAE modes is
mathematically the finite set \(\{\mathsf{Mode2},\mathsf{Mode3},\mathsf{Mode5}\}\). \\
Distribution-support types &
Carrier sets for random coins, signing entropy tokens, sample pairs,
hyperball samples, and rejection-sampling states.  These are the domains on
which the sampler laws are defined. \\
Transcript and event types &
Carrier sets for transcript records, counted ROM events, programming sites, and
bad-event indicators.  These are state spaces for hybrid-game invariants. \\
Hardness-game types &
Carrier sets for Module-SIS instances, MLWE instances, bimodal self-target
instances, solutions, and solver outputs.  These are the mathematical domains of
the assumed hard problems. \\
Hash-support types &
Carrier sets for Keccak bytes, lanes, states, and FIPS 202 states.  These model
bit-level and byte-level hash state spaces. \\
\bottomrule
\end{longtable}

The mathematical significance is that no lemma may mix these sets unless the
EasyCrypt type checker admits the expression.  The type checker therefore
enforces, for example, that a sampler seed, a ROM query, a challenge, and a
signature transcript cannot be confused in a reduction.

\section{Operations as functions, predicates, distributions, and bounds}
\label{app:ops-as-functions}

An EasyCrypt operation declaration of the form
\begin{center}
\ecsym{op f : tau1 -> ... -> taun -> tau}
\end{center}
is interpreted as a total mathematical function
\[
  \llbracket f \rrbracket :
  \llbracket \tau_1 \rrbracket \times \cdots \times
  \llbracket \tau_n \rrbracket
  \to \llbracket \tau \rrbracket .
\]
When \(\tau=\mathsf{Bool}\), the operation is a predicate or event.  When
\(\tau=\mathsf{Distr}(\sigma)\), it is a probability law or sampling
distribution.  When \(\tau=\mathbb{R}\), it is a concrete loss, advantage, or
bound term.  When \(\tau\) is a product or record-like tuple, it is a projection
or constructor.

Every operation in the proof surface falls into one of the following
mathematical classes.

\begin{longtable}{p{0.24\textwidth}p{0.66\textwidth}}
\caption{Mathematical reading of all manifest operation declarations.}
\label{tab:op-correspondence}\\
\toprule
Operation class & Mathematical correspondence \\
\midrule
\endfirsthead
\toprule
Operation class & Mathematical correspondence \\
\midrule
\endhead
Constructors and projections &
Functions building or extracting components of keys, signatures, transcripts,
programming sites, random-oracle outputs, sample pairs, and rejection states.
Mathematically these are product constructors and coordinate projections. \\
Well-formedness predicates &
Predicates \(P(x)\) expressing membership in a syntactic or algebraic domain:
correct list length, coefficient range, sparse challenge form, valid vector
dimension, valid transcript field relation, or valid hardness instance. \\
Algebraic operations &
Total functions implementing modular arithmetic, polynomial addition,
negation, subtraction, multiplication, vector arithmetic, matrix-vector
multiplication, high/low-bit extraction, commitment reconstruction, and public
verification equations. \\
Hash-query constructors &
Injective or tagged constructors for message-hash, challenge-hash,
matrix-expansion, and sampler-expansion queries.  Mathematically they partition
the ROM domain into the query forms used by the security proof. \\
Distributions &
Probability laws for bytes, seeds, CRHs, signing randomness, coefficients,
polynomials, vectors, matrices, challenges, structural samples,
bounded-unit samples, checked-hyperball samples, exact-hyperball samples, and
ROM outputs. \\
Events and bad predicates &
Boolean predicates defining rejection, verification success, forgery success,
prequery, collision, reprogramming failure, min-entropy failure, counted ROM
badness, sampler badness, and transcript readiness.  These are events in the
probability spaces of the games. \\
Budgets and loss terms &
Real-valued functions and constants representing query budgets, support-size
lower bounds, ROM collision terms, prequery terms, min-entropy terms,
programming losses, rejection-sampling losses, and final EUF-CMA bounds. \\
Hardness relations &
Predicates and functions defining Module-SIS, MLWE, and bimodal self-target
relations, together with adapter maps from one relation or game to another. \\
Hash-state functions &
Bit and byte functions for Keccak lanes, state indexing, round transformations,
FIPS 202 absorb/squeeze, and SHAKE256 sizing.  These are deterministic functions
over finite bit/byte sequences. \\
\bottomrule
\end{longtable}

This is the mathematical reason the proof can be read as a sequence of
ordinary definitions.  For example, \ecsym{sampler\_expand\_query} is not a
tactic artifact; it is a function from signing randomness to the ROM domain.
The event \ecsym{sampler\_bad\_prequery} is the corresponding predicate that
this query already lies in the joint log containing adversary hash queries and
earlier sampler-expansion self-queries.  The tight bound is then a
finite-support probability statement about a freshly sampled argument to that
function colliding with the combined log.

\section{Modules as games, algorithms, oracles, and reductions}
\label{app:modules-as-games}

An EasyCrypt module is interpreted as a probabilistic imperative algorithm with
local state.  A module type is an interface specifying the procedures available
to such an algorithm.  A functorial module, such as a game parameterized by an
adversary and an oracle, is interpreted as a mathematical construction sending
algorithms to algorithms:
\[
  (A,H,O) \longmapsto G^{A,H,O}.
\]
The procedures of the resulting module define random variables and events by
executing probabilistic code.

\begin{longtable}{p{0.24\textwidth}p{0.66\textwidth}}
\caption{Mathematical reading of all manifest module declarations.}
\label{tab:module-correspondence}\\
\toprule
Module class & Mathematical correspondence \\
\midrule
\endfirsthead
\toprule
Module class & Mathematical correspondence \\
\midrule
\endhead
Oracle interfaces &
Abstract access to random-oracle procedures and programmable-oracle procedures.
Mathematically these are black-box functions or stateful lazy-sampling oracles. \\
Scheme interfaces &
Abstract key-generation, signing, and verification algorithms.  These define the
cryptographic primitive whose correctness and security games are studied. \\
Adversary interfaces &
Classes of algorithms with specified oracle access.  EUF-CMA adversaries may
call both hash and signing oracles; UF-NMA adversaries receive only the public
interface needed by the no-message game. \\
Security games &
Probabilistic experiments for correctness, EUF-CMA, SUF-CMA, UF-NMA,
ROM-internal transcript games, simulated-signing games, and exact-hyperball
games.  Their outputs are Bernoulli random variables representing success
events. \\
Adapters &
Program transformations from one adversary or solver interface to another, such
as NMA-as-CMA adapters and bimodal-to-Module-SIS adapters.  Mathematically these
are reductions between experiments. \\
Lazy ROM modules &
Stateful samplers implementing the random oracle by memoization.  Their state is
a finite map from queries to sampled outputs plus logs and bad-event flags. \\
Programming modules &
Algorithms that reprogram or observe ROM sites.  Mathematically these implement
the standard ROM programming operation subject to freshness and min-entropy
preconditions. \\
Signing simulators and samplers &
Algorithms that generate signatures or transcript samples under real,
structural, paper, rejection, RO-attempt, or exact-hyperball distributions.
Mathematically these are Markov kernels between message/context inputs and
signature/transcript outputs. \\
Hardness games and reductions &
Experiments defining the abstract success probability of MLWE, Module-SIS, and
bimodal self-target MSIS solvers, plus reductions translating HAETAE forgeries
into solver successes. \\
\bottomrule
\end{longtable}

This module interpretation is essential because the proof is not merely a
collection of static equations.  EUF-CMA security is a statement about an
interactive experiment.  Modules are the formal objects whose executions define
that experiment.

\section{Lemmas as propositions, invariants, equivalences, and inequalities}
\label{app:lemmas-as-theorems}

An EasyCrypt lemma declaration of the form
\begin{center}
\ecsym{lemma L : Phi.}
\end{center}
is interpreted as a theorem asserting the formula
\(\llbracket \Phi \rrbracket\).  The proof script following the lemma is a
machine-checked derivation in EasyCrypt's logic and program logic.  For the
HAETAE security proof, every lemma belongs to one of the following mathematical
families.

\begin{longtable}{p{0.24\textwidth}p{0.66\textwidth}}
\caption{Mathematical reading of all manifest lemma declarations.}
\label{tab:lemma-correspondence}\\
\toprule
Lemma family & Mathematical correspondence \\
\midrule
\endfirsthead
\toprule
Lemma family & Mathematical correspondence \\
\midrule
\endhead
Parameter lemmas &
Arithmetic facts about constants and mode-dependent parameters, such as
positivity, size lower bounds, and challenge-weight feasibility. \\
Algebraic lemmas &
Equational and well-formedness theorems for coefficients, polynomials, vectors,
matrices, challenges, signatures, and verification equations. \\
Distribution lemmas &
Losslessness, support, point-probability, boundedness, and sampler-correctness
facts for the distributions used in signing and simulation. \\
Transcript lemmas &
Theorems saying that transcript predicates imply the public equations,
challenge-query identities, field consistency, and extraction conditions needed
by the hardness reductions. \\
Event lemmas &
Implications relating rejection-free events, acceptance events, verification
success, and forgery events. \\
ROM lemmas &
Losslessness and projection facts for query-dependent ROM output distributions,
plus preservation and counting facts for lazy ROM state. \\
ROM-programming lemmas &
Freshness, non-prequery, min-entropy, counted-event, and reprogramming-failure
bounds.  Mathematically these are fundamental-lemma and union-bound steps over
finite query logs. \\
Invariant lemmas &
State preservation theorems for transcript logs, internal transcript state,
sampler ROM coverage, signing budgets, active public keys, and adversary query
logs. \\
Game-hop lemmas &
Inequalities or equalities between success probabilities of adjacent games.
These are the checked hybrid steps of the security proof. \\
Reduction lemmas &
Theorems transforming a successful forgery or bimodal event into success in an
abstract hardness game, with the corresponding advantage inequality. \\
Concrete-bound lemmas &
Real-arithmetic theorems normalizing, grouping, proving non-negativity of, or
bounding the final concrete loss expression. \\
Top-level security lemmas &
Compositions of all previous families into correctness and conditional EUF-CMA
advantage theorems for the HAETAE model. \\
\bottomrule
\end{longtable}

In mathematician's language, the 1088 checked lemmas form a dependency graph of
propositions.  The leaves are parameter, algebraic, distributional, and
well-formedness facts.  The middle nodes are invariant and hybrid-game
propositions.  The root is the conditional EUF-CMA security implication.
Machine checking means that every edge in this graph is justified by
EasyCrypt's proof kernel, not by an informal appeal to a paper proof.

\section{Hoare and equivalence judgments as program-logic theorems}
\label{app:hoare-equiv-as-program-logic}

The Hoare and equivalence commands are the formal counterpart of reasoning about
program executions.  If \(c\) is a probabilistic procedure, an EasyCrypt Hoare
judgment \ecsym{hoare [c : P ==> Q]}
is interpreted as partial-correctness reasoning: every terminating execution
starting in a state satisfying \(P\) ends in a state satisfying \(Q\).  A
probabilistic Hoare judgment \ecsym{phoare [c : P ==> Q] = p}
is interpreted as a statement about the probability that \(Q\) holds after
executing \(c\) from a \(P\)-state.  An equivalence judgment
\ecsym{equiv [c1 ~ c2 : Rpre ==> Rpost]}
is interpreted as a coupling or relational program theorem: two executions can
be related so that the post-relation holds whenever the pre-relation holds.
The command \ecsym{byequiv} is the proof-script form of such a relational
argument.

The 105 Hoare/equivalence proof commands in the manifest have the following
mathematical roles.

\begin{longtable}{p{0.26\textwidth}p{0.64\textwidth}}
\caption{Mathematical reading of all manifest Hoare/equivalence proof commands.}
\label{tab:hoare-correspondence}\\
\toprule
Program-logic family & Mathematical correspondence \\
\midrule
\endfirsthead
\toprule
Program-logic family & Mathematical correspondence \\
\midrule
\endhead
Correctness Hoare proofs &
Show that modeled key generation, signing, and verification preserve the
correctness predicate and end in accepting states. \\
Adapter equivalences &
Show that wrapping one adversary or solver interface into another does not
change the relevant success event.  Mathematically these are exact reductions
with zero loss. \\
Simulated-signing equivalences &
Show that replacing real signing by transcript or internal-transcript signing
preserves the adversary's observable behavior up to the stated bad event. \\
ROM-internal equivalences &
Show that moving hash behavior into a lazy internal ROM preserves query/output
relations and adversary result distributions under synchronized states. \\
Sampler equivalences &
Show that structural, paper, rejection, RO-attempt, and exact-hyperball samplers
produce related outputs under the coupling conditions proved by sampler lemmas. \\
Fundamental-lemma proofs &
Show that two games are identical on the complement of a bad event, so their
success probabilities differ by at most the probability of that bad event. \\
Budgeted oracle-call proofs &
Show that the budgeted wrappers preserve capped counters and logs, and that the
calls recorded before the cap are exactly the calls counted in the loss term.
Over-budget calls are handled by the modeled wrapper behavior rather than by a
syntactic restriction on arbitrary adversary code. \\
\bottomrule
\end{longtable}

These judgments are the place where informal ``the games are identical until
bad'' reasoning becomes mathematical.  The proof must name the program states,
state relation, bad event, and postcondition.  EasyCrypt then checks that every
procedure call, random sample, assignment, branch, and oracle invocation
preserves the required relation or incurs the stated loss.

\section{File-wise mathematical correspondence}
\label{app:file-wise-correspondence}

The previous sections describe the universal interpretation of the EasyCrypt
object classes.  The following table records how those classes specialize in
each manifest file.

\begin{longtable}{p{0.20\textwidth}p{0.70\textwidth}}
\caption{Per-file mathematical correspondence of proof-object classes.}
\label{tab:file-wise-correspondence}\\
\toprule
File & Mathematical role of its types, operations, modules, lemmas, and
program-logic judgments \\
\midrule
\endfirsthead
\toprule
File & Mathematical role of its types, operations, modules, lemmas, and
program-logic judgments \\
\midrule
\endhead
\ecfile{Sig\_ROM.ec} &
Types are abstract carrier sets for a generic ROM signature scheme.  Operations
are freshness predicates and ROM-output distributions.  Modules are correctness,
EUF-CMA, SUF-CMA, UF-NMA games and adversary/oracle interfaces.  Lemmas prove
freshness base cases and the exact NMA-as-CMA correspondence.  Equivalence
judgments are exact game-isomorphism statements. \\
\ecfile{HAETAE\_Params.ec} &
The mode type is a finite parameter set.  Operations are numeric parameter
functions.  Lemmas are arithmetic facts ensuring positivity, dimensions, and
feasibility side conditions.  There are no stateful games here. \\
\ecfile{HAETAE\_Algebra.ec} &
Types are algebraic carrier sets for HAETAE objects.  Operations are modular,
polynomial, vector, matrix, packing, challenge, and verification functions.
Lemmas are algebraic equations and well-formedness preservation theorems. \\
\ecfile{HAETAE\_Keccak1600.ec} &
Types are bit/byte/lane/state carrier sets.  Operations are Keccak state
transformations and projections.  Lemmas are bit-level equations,
well-formedness facts, and size facts for the permutation model. \\
\ecfile{HAETAE\_FIPS202.ec} &
Types are FIPS 202 byte and state spaces.  Operations are SHAKE256
absorb/squeeze functions and constants.  Lemmas are state-size and output-size
theorems supporting the modeled hash layer. \\
\ecfile{HAETAE\_Distributions.ec} &
Types are sampler domains and sample-pair spaces.  Operations are probability
laws, support predicates, point-bound expressions, and sample constructors.
Lemmas are losslessness, support, boundedness, and sampler-correctness
theorems. \\
\ecfile{HAETAE\_ROM.ec} &
Types are HAETAE ROM query and output spaces.  Operations are tagged query
constructors, output projections, output injections, and query-dependent output
laws.  Modules instantiate the full ROM.  Lemmas prove losslessness and sampler
projection laws. \\
\ecfile{HAETAE\_Scheme.ec} &
The module is the HAETAE instance of the generic signature interface.  The mode
operation fixes the parameter family.  Mathematically this file identifies the
generic game carrier sets with the HAETAE carrier sets. \\
\ecfile{HAETAE\_Events.ec} &
Operations are named events for rejection-freeness and verification acceptance.
Lemmas show that the current modeled events imply the event predicates consumed
by correctness and game-hop theorems. \\
\ecfile{HAETAE\_Transcript.ec} &
The transcript type is the carrier set of signing transcripts.  Operations are
field projections, consistency predicates, validity predicates, forking
relations, and extraction maps.  Lemmas prove that valid transcripts encode the
algebraic equations required by the hardness reductions. \\
\ecfile{HAETAE\_Rejection.ec} &
Types are signing-attempt states and sample spaces.  Operations are rejection
predicates, state projections, signature constructors, acceptance predicates,
and sampler predicates.  Lemmas prove verification, simulation,
sampler-correctness, and rejection-loss facts. \\
\ecfile{HAETAE\_Assumptions.ec} &
Types are hardness-instance and solution spaces.  Operations are relation
predicates and distribution adapters.  Modules are hardness games, solver
interfaces, and reductions.  Lemmas and equivalence judgments prove exactness of
adapters and well-formedness of lifted instances. \\
\ecfile{HAETAE\_Reductions.ec} &
Operations are real-valued hardness terms, query budgets, support bounds, ROM
losses, rejection losses, and final advantage bounds.  Lemmas are arithmetic
normalization, non-negativity, subunit, grouping, and support-dominance
theorems. \\
\ecfile{HAETAE\_ROM\_Programming.ec} &
Types are programming sites and counted ROM events.  Operations are event
counters, freshness predicates, min-entropy predicates, and counted bad events.
Modules are lazy ROM and programming games.  Lemmas are preservation,
freshness, count, min-entropy, and ROM-programming loss theorems. \\
\ecfile{HAETAE\_HopGames.ec} &
Types and operations define transcript invariants, sampler coverage, public
simulation fields, paper-simulation samples, bad events, and budget disciplines.
Modules are the full sequence of real, simulated, transcript, ROM-internal,
NMA, paper-sampler, RO-attempt, exact-hyperball, and budgeted-wrapper games.
Lemmas and Hoare/equivalence judgments prove the hybrid steps, invariants,
fundamental-lemma bounds, sampler-freshness bounds, and budgeted lifting
theorems. \\
\ecfile{HAETAE\_Security.ec} &
Operations define the top-level correctness obligation.  Lemmas compose all
lower-level facts into correctness and conditional EUF-CMA theorems.  Its Hoare
proof is the final correctness program-logic obligation.  Mathematically this
file is the root of the proof DAG: it turns local game-hop and arithmetic
theorems, under the explicit hop and hardness premises of the final composition
lemma, into the published security inequality. \\
\bottomrule
\end{longtable}

\section{Mathematical reading of the final proof DAG}
\label{app:final-proof-dag-reading}

The entire EasyCrypt development can be summarized as the following mathematical
dependency chain.

\[
\begin{array}{c}
\text{parameters and algebraic carrier sets} \\
\Downarrow \\
\text{well-formed distributions and samplers} \\
\Downarrow \\
\text{scheme, transcript, rejection, and ROM events} \\
\Downarrow \\
\text{stateful games and lazy-ROM/programming modules} \\
\Downarrow \\
\text{Hoare/equivalence proofs for adjacent hybrids} \\
\Downarrow \\
\text{bad-event probability bounds and hardness reductions} \\
\Downarrow \\
\text{concrete arithmetic normalization} \\
\Downarrow \\
\text{HAETAE correctness and EUF-CMA security theorem.}
\end{array}
\]

In this reading, Types supply the spaces, Ops supply the maps and events,
Modules supply the experiments, Lemmas supply the propositions, and
Hoare/equiv judgments supply the program-logic bridges between experiments.
The machine checker verifies that these objects compose without type errors,
unproved proof holes, or missing side conditions, thereby converting the
pen-and-paper game proof into a formal theorem about the modeled HAETAE scheme.
